% GENERATED FILE. Edit canonical lessons, then run npm run generate:books.
% canonical-source-hash: fnv1a64:6a03eed7fe98a5bd
% canonical-lessons: PA-C266-mahalla, PA-C266-kasba, PA-C266-hatti, PA-C266-dhaba, PA-C266-imarat, PA-R266-review-260-266

\chapter{Around the Town}
\label{ch:pa-a2-around-the-town}
\hlchaptermodality{266}

\begin{chapteropening}
\textbf{By the end of this chapter:} \emph{I can name the places in a small town.}

Complete the last lesson of chapter 266: I can name the places in a small town.
\end{chapteropening}

\section[\texorpdfstring{\pa{ਮਹੱਲਾ} (mahallā)}{ਮਹੱਲਾ (mahallā)}]{\pa{ਮਹੱਲਾ} (mahallā) — a neighbourhood; a quarter of a town}

\label{lesson:PA-C266-mahalla}



\begin{quote}
Before the new one: say the Punjabi for the distance to travel, then the Punjabi for a kos.
\end{quote}

\subsection*{You'll want to know: \pa{ਮਹੱਲਾ}}
\textbf{\pa{ਮਹੱਲਾ}} — \emph{mahallā} — \textquotedblleft{}a neighbourhood; a quarter of a town\textquotedblright{}.

\begin{quote}
\pa{ਸਾਡਾ} \pa{ਮਹੱਲਾ} \pa{ਬਹੁਤ} \pa{ਸ਼ਾਂਤ} \pa{ਹੈ।}
\end{quote}

\begin{quote}
\emph{sāḍā mahallā bahut shā̃t hai.}
\end{quote}

\begin{quote}
Our neighbourhood is very calm.
\end{quote}

A masculine noun.

\begin{cousinweb}
\textbf{\pa{ਮਹੱਲਾ}} (\emph{mahallā}) is a loan, usually traced through Persian to Arabic \emph{mahalla}, \textquotedblleft{}a quarter of a town\textquotedblright{}.
\end{cousinweb}

\begin{grammarlens}[title={What you've built}]
One more word for this chapter.
\end{grammarlens}

\subsection*{Your turn}
\begin{itemize}
\raggedright
  \item \emph{Say it:} \emph{mahallā}
  \item \emph{Say it:} \emph{mahallā}, once more
  \item \emph{Say it:} \emph{sāḍā mahallā bahut shā̃t hai.} — Our neighbourhood is very calm.
  \item \emph{Recall:} say the Punjabi for the distance to travel, then the Punjabi for a kos, then say \emph{mahallā} again
\end{itemize}

\begin{tcolorbox}[breakable,colback=teal!4,colframe=teal!35!black,title={Before you move on}]
What does \pa{ਮਹੱਲਾ} mean? (\textquotedblleft{}A neighbourhood; a quarter of a town\textquotedblright{}.) Say it once more.
\end{tcolorbox}
\section[\texorpdfstring{\pa{ਕਸਬਾ} (kasbā)}{ਕਸਬਾ (kasbā)}]{\pa{ਕਸਬਾ} (kasbā) — a small town}

\label{lesson:PA-C266-kasba}



\begin{quote}
Before the new one: say the Punjabi for a kos, then the Punjabi for a neighbourhood.
\end{quote}

\subsection*{You'll want to know: \pa{ਕਸਬਾ}}
\textbf{\pa{ਕਸਬਾ}} — \emph{kasbā} — \textquotedblleft{}a small town\textquotedblright{}.

\begin{quote}
\pa{ਮੇਰਾ} \pa{ਦੋਸਤ} \pa{ਇੱਕ} \pa{ਛੋਟੇ} \pa{ਕਸਬੇ} \pa{ਵਿੱਚ} \pa{ਰਹਿੰਦਾ} \pa{ਹੈ।}
\end{quote}

\begin{quote}
\emph{merā dost ikk chhoṭe kasbe vich rahindā hai.}
\end{quote}

\begin{quote}
My friend lives in a small town.
\end{quote}

A masculine noun.

\begin{cousinweb}
\textbf{\pa{ਕਸਬਾ}} (\emph{kasbā}) is a loan, usually traced through Persian to Arabic \emph{qasaba}, \textquotedblleft{}a town\textquotedblright{}.
\end{cousinweb}

\begin{grammarlens}[title={What you've built}]
One more word for this chapter.
\end{grammarlens}

\subsection*{Your turn}
\begin{itemize}
\raggedright
  \item \emph{Say it:} \emph{kasbā}
  \item \emph{Say it:} \emph{kasbā}, once more
  \item \emph{Say it:} \emph{merā dost ikk chhoṭe kasbe vich rahindā hai.} — My friend lives in a small town.
  \item \emph{Recall:} say the Punjabi for a kos, then the Punjabi for a neighbourhood, then say \emph{kasbā} again
  \item \emph{Write it:} copy \textbf{\pa{ਕਸਬਾ}} once, checking the whole word
\end{itemize}

\begin{tcolorbox}[breakable,colback=teal!4,colframe=teal!35!black,title={Before you move on}]
What does \pa{ਕਸਬਾ} mean? (\textquotedblleft{}A small town\textquotedblright{}.) Say it once more.
\end{tcolorbox}
\section[\texorpdfstring{\pa{ਹੱਟੀ} (haṭṭī)}{ਹੱਟੀ (haṭṭī)}]{\pa{ਹੱਟੀ} (haṭṭī) — a small shop}

\label{lesson:PA-C266-hatti}



\begin{quote}
Before the new one: say the Punjabi for a neighbourhood, then the Punjabi for a small town.
\end{quote}

\subsection*{You'll want to know: \pa{ਹੱਟੀ}}
\textbf{\pa{ਹੱਟੀ}} — \emph{haṭṭī} — \textquotedblleft{}a small shop\textquotedblright{}.

\begin{quote}
\pa{ਗਲੀ} \pa{ਵਿੱਚ} \pa{ਇੱਕ} \pa{ਹੱਟੀ} \pa{ਹੈ} \pa{ਅਤੇ} \pa{ਉੱਥੇ} \pa{ਦੁੱਧ} \pa{ਮਿਲਦਾ} \pa{ਹੈ।}
\end{quote}

\begin{quote}
\emph{galī vich ikk haṭṭī hai ate utthe dudh mildā hai.}
\end{quote}

\begin{quote}
There is a little shop in the lane, and you can get milk there.
\end{quote}

A feminine noun.

\begin{cousinweb}
\textbf{\pa{ਹੱਟੀ}} (\emph{haṭṭī}) is inherited from Sanskrit \emph{haṭṭa}, \textquotedblleft{}market, fair\textquotedblright{}, beside \emph{haṭṭī}, \textquotedblleft{}petty market\textquotedblright{}; CDIAL lists Punjabi \emph{haṭṭ}, \textquotedblleft{}shop\textquotedblright{}, under it (CDIAL 13944).
\end{cousinweb}

\begin{grammarlens}[title={What you've built}]
One more word for this chapter.
\end{grammarlens}

\subsection*{Your turn}
\begin{itemize}
\raggedright
  \item \emph{Say it:} \emph{haṭṭī}
  \item \emph{Say it:} \emph{haṭṭī}, once more
  \item \emph{Say it:} \emph{galī vich ikk haṭṭī hai ate utthe dudh mildā hai.} — There is a little shop in the lane, and you can get milk there.
  \item \emph{Recall:} say the Punjabi for a neighbourhood, then the Punjabi for a small town, then say \emph{haṭṭī} again
  \item \emph{Write it:} copy \textbf{\pa{ਹੱਟੀ}} once, checking the whole word
\end{itemize}

\begin{tcolorbox}[breakable,colback=teal!4,colframe=teal!35!black,title={Before you move on}]
What does \pa{ਹੱਟੀ} mean? (\textquotedblleft{}A small shop\textquotedblright{}.) Say it once more.
\end{tcolorbox}
\section[\texorpdfstring{\pa{ਢਾਬਾ} (ḍhābā)}{ਢਾਬਾ (ḍhābā)}]{\pa{ਢਾਬਾ} (ḍhābā) — a roadside eatery}

\label{lesson:PA-C266-dhaba}



\begin{quote}
Before the new one: say the Punjabi for a small town, then the Punjabi for a small shop.
\end{quote}

\subsection*{You'll want to know: \pa{ਢਾਬਾ}}
\textbf{\pa{ਢਾਬਾ}} — \emph{ḍhābā} — \textquotedblleft{}a roadside eatery\textquotedblright{}.

\begin{quote}
\pa{ਅਸੀਂ} \pa{ਸੜਕ} \pa{ਦੇ} \pa{ਨੇੜੇ} \pa{ਢਾਬੇ} \pa{ਉੱਤੇ} \pa{ਖਾਣਾ} \pa{ਖਾਂਦੇ} \pa{ਹਾਂ।}
\end{quote}

\begin{quote}
\emph{asīṁ saṛak de neṛe ḍhābe utte khāṇā khānde hāṁ.}
\end{quote}

\begin{quote}
We eat at a roadside eatery near the road.
\end{quote}

A masculine noun.

\begin{cousinweb}
\textbf{\pa{ਢਾਬਾ}} (\emph{ḍhābā}) has no settled source: the Punjabi origin remains unknown in the sources used here.
\end{cousinweb}

\begin{grammarlens}[title={What you've built}]
One more word for this chapter.
\end{grammarlens}

\subsection*{Your turn}
\begin{itemize}
\raggedright
  \item \emph{Say it:} \emph{ḍhābā}
  \item \emph{Say it:} \emph{ḍhābā}, once more
  \item \emph{Say it:} \emph{asīṁ saṛak de neṛe ḍhābe utte khāṇā khānde hāṁ.} — We eat at a roadside eatery near the road.
  \item \emph{Recall:} say the Punjabi for a small town, then the Punjabi for a small shop, then say \emph{ḍhābā} again
  \item \emph{Write it:} copy \textbf{\pa{ਢਾਬਾ}} once, checking the whole word
\end{itemize}

\begin{tcolorbox}[breakable,colback=teal!4,colframe=teal!35!black,title={Before you move on}]
What does \pa{ਢਾਬਾ} mean? (\textquotedblleft{}A roadside eatery\textquotedblright{}.) Say it once more.
\end{tcolorbox}
\section[\texorpdfstring{\pa{ਇਮਾਰਤ} (imārat)}{ਇਮਾਰਤ (imārat)}]{\pa{ਇਮਾਰਤ} (imārat) — a building}

\label{lesson:PA-C266-imarat}



\begin{quote}
Before the new one: say the Punjabi for a small shop, then the Punjabi for a roadside eatery.
\end{quote}

\subsection*{You'll want to know: \pa{ਇਮਾਰਤ}}
\textbf{\pa{ਇਮਾਰਤ}} — \emph{imārat} — \textquotedblleft{}a building\textquotedblright{}.

\begin{quote}
\pa{ਬੈਂਕ} \pa{ਇੱਕ} \pa{ਵੱਡੀ} \pa{ਇਮਾਰਤ} \pa{ਵਿੱਚ} \pa{ਹੈ।}
\end{quote}

\begin{quote}
\emph{baiṅk ikk vaḍḍī imārat vich hai.}
\end{quote}

\begin{quote}
The bank is in a big building.
\end{quote}

A feminine noun.

\begin{cousinweb}
\textbf{\pa{ਇਮਾਰਤ}} (\emph{imārat}) is a loan, usually traced through Persian to Arabic \emph{imāra}, \textquotedblleft{}a building\textquotedblright{}.
\end{cousinweb}

\begin{grammarlens}[title={What you've built}]
One more word for this chapter.
\end{grammarlens}

\subsection*{Your turn}
\begin{itemize}
\raggedright
  \item \emph{Say it:} \emph{imārat}
  \item \emph{Say it:} \emph{imārat}, once more
  \item \emph{Say it:} \emph{baiṅk ikk vaḍḍī imārat vich hai.} — The bank is in a big building.
  \item \emph{Recall:} say the Punjabi for a small shop, then the Punjabi for a roadside eatery, then say \emph{imārat} again
  \item \emph{Write it:} copy \textbf{\pa{ਇਮਾਰਤ}} once, checking the whole word
\end{itemize}

\begin{tcolorbox}[breakable,colback=teal!4,colframe=teal!35!black,title={Before you move on}]
What does \pa{ਇਮਾਰਤ} mean? (\textquotedblleft{}A building\textquotedblright{}.) Say it once more.
\end{tcolorbox}
\section[(dialogue)]{Looking back: from Quarter Past and Half Past to Around the Town}

\label{lesson:PA-R266-review-260-266}



\begin{quote}
Say the Punjabi for plus a quarter and plus a half.
\end{quote}

\subsection*{Your turn}
Say these aloud:

\begin{itemize}
\raggedright
  \item plus a quarter, plus a half, less a quarter, one and a half, two and a half
  \item half, a railway platform, to halt, to set off, departing
  \item a fee, a fare, a fine, a receipt, to pay off
  \item a bend in the road, to turn, a side, right-hand, left-hand
  \item across, to pass through, a big gate, a pole, a corner
  \item on this side, farther off, in the middle, the distance to travel, a kos
  \item a neighbourhood, a small town, a small shop, a roadside eatery, a building
\end{itemize}

\begin{tcolorbox}[breakable,colback=teal!4,colframe=teal!35!black,title={Before you move on}]
Plus a quarter? (\textbf{\pa{ਸਵਾ}}, \emph{savā}.) Plus a half? (\textbf{\pa{ਸਾਢੇ}}, \emph{sāḍhe}.) Less a quarter? (\textbf{\pa{ਪੌਣੇ}}, \emph{pauṇe}.) One and a half? (\textbf{\pa{ਡੇਢ}}, \emph{ḍeḍh}.) Two and a half? (\textbf{\pa{ਢਾਈ}}, \emph{ḍhāī}.) Half? (\textbf{\pa{ਅੱਧਾ}}, \emph{addhā}.) A railway platform? (\textbf{\pa{ਪਲੇਟਫ਼ਾਰਮ}}, \emph{paleṭfāram}.) To halt? (\textbf{\pa{ਠਹਿਰਨਾ}}, \emph{ṭhahirnā}.) To set off? (\textbf{\pa{ਚੱਲਣਾ}}, \emph{challṇā}.) Departing? (\textbf{\pa{ਰਵਾਨਾ}}, \emph{ravānā}.) A fee? (\textbf{\pa{ਫ਼ੀਸ}}, \emph{fīs}.) A fare? (\textbf{\pa{ਕਿਰਾਇਆ}}, \emph{kirāiā}.) A fine? (\textbf{\pa{ਜੁਰਮਾਨਾ}}, \emph{jurmānā}.) A receipt? (\textbf{\pa{ਰਸੀਦ}}, \emph{rasīd}.) To pay off? (\textbf{\pa{ਚੁਕਾਉਣਾ}}, \emph{chukāuṇā}.) A bend in the road? (\textbf{\pa{ਮੋੜ}}, \emph{moṛ}.) To turn? (\textbf{\pa{ਮੁੜਨਾ}}, \emph{muṛnā}.) A side? (\textbf{\pa{ਪਾਸਾ}}, \emph{pāsā}.) Right-hand? (\textbf{\pa{ਸੱਜਾ}}, \emph{sajjā}.) Left-hand? (\textbf{\pa{ਖੱਬਾ}}, \emph{khabbā}.) Across? (\textbf{\pa{ਪਾਰ}}, \emph{pār}.) To pass through? (\textbf{\pa{ਲੰਘਣਾ}}, \emph{laṅghṇā}.) A big gate? (\textbf{\pa{ਫਾਟਕ}}, \emph{phāṭak}.) A pole? (\textbf{\pa{ਖੰਭਾ}}, \emph{khambhā}.) A corner? (\textbf{\pa{ਕੋਨਾ}}, \emph{konā}.) On this side? (\textbf{\pa{ਉਰੇ}}, \emph{ure}.) Farther off? (\textbf{\pa{ਪਰੇ}}, \emph{pare}.) In the middle? (\textbf{\pa{ਵਿਚਕਾਰ}}, \emph{vichkār}.) The distance to travel? (\textbf{\pa{ਵਾਟ}}, \emph{vāṭ}.) A kos? (\textbf{\pa{ਕੋਹ}}, \emph{koh}.) A neighbourhood? (\textbf{\pa{ਮਹੱਲਾ}}, \emph{mahallā}.) A small town? (\textbf{\pa{ਕਸਬਾ}}, \emph{kasbā}.) A small shop? (\textbf{\pa{ਹੱਟੀ}}, \emph{haṭṭī}.) A roadside eatery? (\textbf{\pa{ਢਾਬਾ}}, \emph{ḍhābā}.) A building? (\textbf{\pa{ਇਮਾਰਤ}}, \emph{imārat}.)
\end{tcolorbox}
